% TOP_PROBLEM: {"id": 2304021, "title": "Research Problems in Function Theory — Problem 4.21", "category_id": 8, "category": {"id": 8, "name": "analysis", "display_name": "Analysis", "description": "Limits, continuity, calculus, and function theory.", "slug": "analysis", "order_index": 8, "created_at": "2026-07-31T15:26:25.671Z"}, "status": "open"}
% FIRST_POSTED: null
% ENTRY_KIND: statement_only
% Imported from: batch13.tex; UnsolvedMath ID 2304021
\documentclass[11pt]{book}
\usepackage{fontspec}
\setmainfont{Latin Modern Roman}
\setsansfont{Latin Modern Sans}
\usepackage{amsmath,amssymb,mathtools}
\usepackage{unicode-math}
\setmathfont{Latin Modern Math}
\usepackage{xurl}
\usepackage[hidelinks]{hyperref}
\newcommand{\sref}[2]{\hyperlink{source:#1:#2}{[S#2]}}
\newcommand{\cataloguescope}[1]{\paragraph{Mathematical question.}#1}
\begin{document}
% UnsolvedMath ID 2304021; source catalogue code AMR-022-4021
\setcounter{section}{2304020}
\section{Research Problems in Function Theory — Problem 4.21}\label{2304021}
{\small\sffamily UnsolvedMath ID 2304021 · Analysis}
\subsection{Definitions and mathematical statement}

\paragraph{Shared notation.}$\mathbb N=\{1,2,\ldots\}$, and $\mathbb Z,\mathbb Q,\mathbb R,\mathbb C$ denote the integers, rationals, reals and complex numbers. Any other conventions are specified locally.


Let $n\ge1$ and $a_0,\ldots,a_n\in\{-1,1\}$. For $1\le k\le n$, define the aperiodic correlations
\[b_k=\sum_{j=0}^{n-k}a_{j+k}a_j.
\]
Put $p(z)=\sum_{j=0}^na_jz^j$. Direct expansion and Fourier orthogonality give
\[|p(e^{i\theta})|^2=(n+1)+2\sum_{k=1}^nb_k\cos(k\theta),\]
\[\frac1{2\pi}\int_0^{2\pi}|p(e^{i\theta})|^4\,d\theta=(n+1)^2+2\sum_{k=1}^nb_k^2.
\]
Thus the auxiliary constant term is $n+1$, corresponding to the $n+1$ coefficients.
\paragraph{Statement recorded in the supplied source.}
Is there an absolute constant $A>0$ such that
\[\sum_{k=1}^n|b_k|^2>An^2\]
for every $n\ge1$ and every sign sequence? Such a bound would in particular imply the source's fourth-mean lower bound $\frac1{2\pi}\int_0^{2\pi}|p(e^{i\theta})|^4\,d\theta\ge n^2(1+A)$. \sref{2304021}{2}
\subsection{Short English statement}
Must the sum of squared aperiodic correlations of every sign sequence be at least a fixed positive multiple of the square of its length parameter?
\subsection{Sources}
\begin{description}
\item[\hypertarget{source:2304021:1}{S1}] ULAM.AI and the UnsolvedMath Contributors, \emph{UnsolvedMath: A Curated Collection of Open Mathematics Problems}, record 2304021 (AMR-022-4021). CC BY 4.0. \url{https://huggingface.co/datasets/ulamai/UnsolvedMath}.
\item[\hypertarget{source:2304021:2}{S2}] Walter K. Hayman and Edward F. Lingham, \emph{Research Problems in Function Theory}, arXiv:1809.07200, 2018. Problem 4.21, its complete Update 4.21, and the relevant chapter notation. \url{https://arxiv.org/pdf/1809.07200}.
\end{description}
\label{end:2304021}
\end{document}
